\documentclass[11pt]{article}
\usepackage[margin=1in]{geometry}
\usepackage{amsmath,amssymb,amsthm,mathtools,microtype,hyperref}
\hypersetup{hidelinks}
\newtheorem{theorem}{Theorem}
\newtheorem{lemma}[theorem]{Lemma}
\newtheorem{corollary}[theorem]{Corollary}
\newtheorem{remark}[theorem]{Remark}
\newcommand{\E}{\mathbb E}
\newcommand{\Pp}{\mathbb P}
\newcommand{\R}{\mathbb R}
\newcommand{\ind}{\mathbf 1}
\title{Global Margin Conditioning Amplifies a Gaussian Posterior Mean}
\author{Dmitry Vaintrob \and Kaarel H\"anni \and OpenAI Codex}
\date{October 2026}
\begin{document}
\maketitle
\begin{abstract}
We study a Gaussian-process classification prior conditioned to have a fixed absolute margin at every point of the Boolean cube, followed by one noiseless label observation. For the fixed-depth iterated-arcsine kernels arising at the $\gamma=0$ neural-network scaling, the posterior mean has squared reproducing-kernel Hilbert space norm at least $(1-o(1))v d\log2$, where $v$ is the smallest matrix covariance eigenvalue. Consequently, at depth at least three, a dimension-uniform mean-energy estimate of order $n/\tau_{\alpha d}$ is false already at $n=1$; here $\tau_{\alpha d}$ is the high-degree spectral trace. The proof retains the global-margin conditioning exactly. It identifies an obstruction to transferring an unconditioned Gaussian learning argument, rather than establishing a sample-complexity bound.
\end{abstract}

\section{Model and result}
Let $\mathcal X_d=\{-1,1\}^d$, $M=2^d$, and fix a depth $L\ge2$.
Set
\[
 A(t)=\frac2\pi\arcsin t,\qquad F=A^{\circ(L-1)},\qquad
 K(x,y)=F\!\left(\frac{1+x\cdot y}{d+1}\right),\qquad
 \beta=2^{-(L-1)}.
\]
The covariance matrix $K$ has diagonal one. The Walsh characters diagonalize the normalized integral operator $K/M$; denote its eigenvalue at degree $j$ by $\lambda_j$, and write
\[
 v=\lambda_{\min}(K)>0,\qquad
 \tau_r=\sum_{j>r}\binom dj\lambda_j.
\]
For a vector $h$ its squared RKHS norm is $\|h\|_K^2=h^TK^{-1}h$.
Let $G\sim N(0,K)$ and fix a margin $a>0$, independent of $d$.
The prior is first conditioned on
\[
 \mathcal A_0=\{|G(x)|>a\text{ for every }x\in\mathcal X_d\}.
\]
It is then conditioned on one noiseless label $y\in\{-1,1\}$ at an arbitrary point $z$, that is, on $yG(z)>0$.
Under $\mathcal A_0$ this is equivalent to $yG(z)>a$.
Write $m_G=\E[G\mid\mathcal A_0,yG(z)>a]$.

It is useful to split off independent noise:
\[
 G=H+\sqrt v\,\xi,\qquad H\sim N(0,C),\quad C=K-vI,
 \qquad \xi\sim N(0,I),
\]
with independence. The possibly singular covariance $C$ is interpreted on its range; let $m_H$ be the conditional mean of $H$ under the same event, and put $\|h\|_C^2=h^TC^+h$ on that range.

\begin{theorem}\label{thm:main}
For every fixed $L\ge2$ and $a>0$, uniformly in the observed point and its label,
\[
 \|m_G\|_K^2\ge(1-o(1))v d\log2,
 \qquad \|m_H\|_C^2\ge(1-o(1))v d\log2
 \qquad(d\longrightarrow\infty).
\]
For every fixed $\alpha\in(0,1/2)$,
\[
 v\ge c_L(d\log d)^{-\beta},\qquad
 \tau_{\lfloor\alpha d\rfloor}\ge c_{L,\alpha}d^{-\beta}.
\]
Therefore, if $L\ge3$, both
\[
 \tau_{\lfloor\alpha d\rfloor}\|m_G\|_K^2\longrightarrow\infty,
 \qquad
 \tau_{\lfloor\alpha d\rfloor}\|m_H\|_C^2\longrightarrow\infty.
\]
In particular, no constant independent of $d$ can bound either mean energy by $C n/\tau_{\lfloor\alpha d\rfloor}$ uniformly in the sample size, since it fails at $n=1$.
\end{theorem}

This theorem concerns the actual posterior after global margin conditioning. It does not exclude a bound with an additive dimension-dependent baseline, an amplitude factor, or a restriction to much larger sample sizes. The divergence conclusion is not claimed at $L=2$. The same one-observation result applies when the observed label is generated by parity. It does not itself establish a lower or upper bound for confidently predicting parity at unobserved inputs.

\section{Exact global-margin conditioning}
Write $\Phi$ and $\phi$ for the standard normal distribution and density, and $Q=1-\Phi$. Define
\[
 B_a(t)=\Phi\!\left(\frac{t-a}{\sqrt v}\right)
       +\Phi\!\left(\frac{-t-a}{\sqrt v}\right),
 \qquad
 p_a(t)=\frac{\Phi((t-a)/\sqrt v)}{B_a(t)}.
\]
Conditioning on $H$, the independent noises give
\begin{equation}\label{eq:latent}
 d\mathcal L(H\mid\mathcal A_0)
 =\frac{\prod_x B_a(H(x))}{\Pp(\mathcal A_0)}\,dN(0,C)(H).
\end{equation}
Moreover $\Pp(G(x)>a\mid H,\mathcal A_0)=p_a(H(x))$.
The function $B_a$ is even and increasing on $[0,\infty)$, as follows by differentiating. Also $p_a(-t)=1-p_a(t)$ and $p_a(t)\ge1/2$ for $t\ge0$.

\begin{lemma}[A denominator from a constant shift]\label{lem:denom}
Set $t_0=\sqrt{2\log(4M)}$ and $R=a+t_0+1$. Then
\[
 \Pp(\mathcal A_0)\ge\frac12\exp\!\left(-\frac{R^2}{2\lambda_0}\right).
\]
\end{lemma}
\begin{proof}
The event $B=\{\max_x|G(x)|\le t_0\}$ has probability at least $1/2$ by the Gaussian tail bound and a union bound. The constant vector $h=R\ind$ has energy $R^2/\lambda_0$, and $h+B\subset\mathcal A_0$. Translating the Gaussian density gives
\[
 \Pp(G\in h+B)
 =e^{-\|h\|_K^2/2}\E[\ind_B(G)e^{-h^TK^{-1}G}]
 \ge e^{-\|h\|_K^2/2}\Pp(B).
\]
The last inequality uses the symmetry of $B$ and $\cosh u\ge1$.
\end{proof}

\begin{lemma}[Latent magnitudes forced by all-input conditioning]\label{lem:margin}
Suppose $D\ge1$, $D<M$, and $-\log\Pp(\mathcal A_0)\le D$. Set
\[
 t=\sqrt{\log(M/D)},\quad b=a+\sqrt v\,t,\quad
 c_* =\frac{e^{-3/2}}{\sqrt{2\pi}}.
\]
If $e^{-2a^2/v}\le1/2$, then, for every fixed $z\in\mathcal X_d$,
\begin{equation}\label{eq:q}
 \Pp(|H(z)|\le b\mid\mathcal A_0)
 \le \frac{2}{c_*}\frac{(t+1)(D+1)}{\sqrt{MD}}.
\end{equation}
\end{lemma}
\begin{proof}
Let $q_b(H)=M^{-1}\#\{x:|H(x)|\le b\}$ and $c_b=-\log B_a(b)>0$.
On $\{q_b\ge u\}$, the density numerator in~\eqref{eq:latent} is at most $e^{-M c_bu}$. Thus
\[
 \Pp(q_b\ge u\mid\mathcal A_0)\le\min\{1,e^{D-Mc_bu}\},
 \qquad \E[q_b\mid\mathcal A_0]\le\frac{D+1}{M c_b}.
\]
The expectation bound follows by integrating the tail on $u\ge0$.
For $c\ge0$ and $t\ge0$, shifting the Gaussian integral gives
$Q(t+c)\le e^{-tc-c^2/2}Q(t)$. Therefore
\[
 1-B_a(b)=Q(t)-Q(t+2a/\sqrt v)\ge Q(t)/2.
\]
Integrating $\phi$ from $t$ to $t+(t+1)^{-1}$ yields
$Q(t)\ge c_*e^{-t^2/2}/(t+1)$.
Since $-\log B\ge1-B$,
\[
 M c_b\ge\frac{c_*}{2}\frac{\sqrt{MD}}{t+1}.
\]
Finally $K$, $C$, and $\mathcal A_0$ are invariant under cube translations, which act transitively. Consequently the probability on the left of~\eqref{eq:q} equals $\E[q_b\mid\mathcal A_0]$.
\end{proof}

\begin{lemma}[One label reveals the large latent amplitude]\label{lem:one}
Suppose $b\ge a$, and set $q=\Pp(|H(z)|\le b\mid\mathcal A_0)$. For the positive observed label,
\[
 m_H(z)\ge b\,[1-2p_a(-b)](1-q),\qquad m_G(z)\ge m_H(z).
\]
The corresponding absolute-value bounds hold for a negative label. Also
\[
 p_a(-b)\le\exp\!\left(-\frac{(a+b)^2}{2v}\right).
\]
\end{lemma}
\begin{proof}
Global sign symmetry gives probability $1/2$ to each label, and $\E[H(z)\mid\mathcal A_0]=0$. Hence
\[
 m_H(z)=\E[H(z)\operatorname{sign}G(z)\mid\mathcal A_0]
 =\E[|H(z)|(2p_a(|H(z)|)-1)\mid\mathcal A_0].
\]
The integrand is nonnegative and increasing in $|H(z)|$. To see the needed monotonicity, the odds $p_a(u)/(1-p_a(u))=\Phi((u-a)/\sqrt v)/\Phi((-u-a)/\sqrt v)$ increase with $u$.
Restricting to $|H(z)|>b$ proves the first bound.
For $u\ge0$, the normal Mills ratio $\phi(u)/Q(u)$ is at least $u$, because it is the conditional mean of a standard normal given that it exceeds $u$.
Integrating gives $Q(u)\le\tfrac12e^{-u^2/2}$.
Since $b\ge a$,
\[
 p_a(-b)\le\frac{Q((a+b)/\sqrt v)}{\Phi((b-a)/\sqrt v)}
 \le e^{-(a+b)^2/(2v)}.
\]
Finally, given $H$, $\mathcal A_0$, and $G(z)>a$, the residual noise at $z$ is a one-sided truncated normal with positive mean. Thus
$\E[G(z)\mid H,\mathcal A_0,G(z)>a]\ge H(z)$, which implies $m_G(z)\ge m_H(z)$. Global sign symmetry handles the negative label.
\end{proof}

\section{Elementary spectral estimates}
The estimates needed here follow without a sharp eigenvalue asymptotic.

\begin{lemma}[Coefficient tails]\label{lem:coeff}
Write $F(u)=\sum_{m\ge0}c_m u^m$. Then $c_m\ge0$, $\sum_m c_m=1$, $c_1=(2/\pi)^{L-1}>0$, and
\[
 c_L R^{-\beta}\le\sum_{m\ge R}c_m\le C_L R^{-\beta}
 \qquad(R\ge2).
\]
\end{lemma}
\begin{proof}
The nonnegative power series of arcsine and its compositions gives the first statements. The expansion
$A(1-s)=1-(2\sqrt2/\pi)s^{1/2}+O(s^{3/2})$
implies by induction that
$1-F(1-s)\asymp_L s^\beta$ as $s\downarrow0$.
Let $T(R)=\sum_{m\ge R}c_m$. Since
\[
 1-F(1-1/R)\ge[1-(1-1/R)^R]T(R),
\]
the upper bound follows. Consequently
$\sum_{m<R}m c_m\le\sum_{j<R}T(j)\le C'_L R^{1-\beta}$.
For any fixed small $\epsilon>0$, Bernoulli's inequality gives
\[
 c_L\epsilon^\beta R^{-\beta}
 \le 1-F(1-\epsilon/R)
 \le C'_L\epsilon R^{-\beta}+T(R).
\]
Choosing $\epsilon$ sufficiently small proves the lower bound, since $\beta<1$.
\end{proof}

\begin{lemma}[Bounds for the constant mode, smallest mode, and tail]\label{lem:spectrum}
For fixed $\alpha\in(0,1/2)$ and sufficiently large $d$,
\[
 \lambda_0\ge\frac{(2/\pi)^{L-1}}{d+1},\quad
 c_L(d\log d)^{-\beta}\le v\le C_L d^{-\beta},\quad
 c_{L,\alpha}d^{-\beta}\le\tau_{\lfloor\alpha d\rfloor}
 \le C_{L,\alpha}d^{-\beta}.
\]
\end{lemma}
\begin{proof}
Put $D=d+1$. The degree-$m$ kernel $[(1+x\cdot y)/D]^m$ has normalized Walsh eigenvalues equal to the law of a walk on subsets of $[d]$: at each step choose uniformly from $\{0,1,\dots,d\}$; choice zero leaves the subset fixed, and choice $i\ge1$ toggles $i$. Write $S_m$ for the resulting subset. Thus
\[
 \lambda_j=\sum_m c_m\Pp(S_m=S)\quad(|S|=j),
 \qquad \tau_r=\sum_m c_m\Pp(|S_m|>r).
\]
The $m=1$ term proves the constant-mode bound. Since $|S_m|\le m$, Lemma~\ref{lem:coeff} proves the upper tail bound.
At least $M/2$ characters have degree greater than $\lfloor d/4\rfloor$ for sufficiently large $d$. Since every matrix eigenvalue is at least $v$, their normalized trace is at least $v/2$; the upper tail bound therefore gives $v\le C_L d^{-\beta}$.

For a lower bound on every eigenvalue, Fourier inversion for the walk gives
\[
 \left|M\Pp(S_m=S)-1\right|
 \le\sum_{j=1}^d\binom dj|1-2j/D|^m
 \le(1+e^{-2m/D})^D-1.
\]
For the second inequality, pair $j$ with $D-j$ and use
$\binom dj+\binom d{j-1}=\binom Dj$ and
$|1-2j/D|\le e^{-2\min(j,D-j)/D}$.
If $m\ge D\log(4D)$, the final expression is less than $1/2$.
It follows that every $\lambda_j$ is at least
$\tfrac1{2M}\sum_{m\ge D\log(4D)}c_m$.
Lemma~\ref{lem:coeff} proves the lower bound on $v$.

Finally
\[
 \E|S_m|=\frac d2\left[1-(1-2/D)^m\right].
\]
Choose a constant $C_\alpha$ so that for every $m\ge C_\alpha D$ this expectation is at least $(\alpha+\eta)d$ for some $\eta>0$.
Since $|S_m|\le d$, it follows that
$\Pp(|S_m|>\lfloor\alpha d\rfloor)\ge\eta/(1-\alpha)$.
Summing over these $m$ and applying Lemma~\ref{lem:coeff} proves the lower tail bound.
\end{proof}

\section{Proof of the theorem and a distinction for posterior draws}
\begin{proof}[Proof of Theorem~\ref{thm:main}]
Take
\[
 D=1+\log2+\frac{(a+\sqrt{2\log(4M)}+1)^2}{\lambda_0}.
\]
Lemmas~\ref{lem:denom} and~\ref{lem:spectrum} give an admissible bound $D=O_{L,a}(d^2)$ on $-\log\Pp(\mathcal A_0)$.
Since $v\to0$, Lemma~\ref{lem:margin} applies for large $d$. With
$b=a+\sqrt{v\log(M/D)}$, it gives $q=o(1)$ in Lemma~\ref{lem:one}.
Also $p_a(-b)=o(1)$, so
\[
 |m_H(z)|\ge(1-o(1))b,\qquad |m_G(z)|\ge(1-o(1))b.
\]
The RKHS evaluation inequality, valid also for a singular covariance on its range, gives
\[
 \|m_H\|_C^2\ge\frac{m_H(z)^2}{C(z,z)}
 =\frac{m_H(z)^2}{1-v},\qquad
 \|m_G\|_K^2\ge\frac{m_G(z)^2}{K(z,z)}=m_G(z)^2.
\]
Because $\log(M/D)=d\log2-O(\log d)$, these prove the energy lower bounds.
The remaining estimates follow from Lemma~\ref{lem:spectrum}. In particular
\[
 v d\,\tau_{\lfloor\alpha d\rfloor}
 \ge c_{L,\alpha}\frac{d^{1-2\beta}}{(\log d)^\beta}
 \longrightarrow\infty\qquad(\beta<1/2).
\]
\end{proof}

A proposed bound for a posterior \emph{draw} has a further, elementary obstruction. For any positive definite $M\times M$ covariance, any positive global margin, and any compatible noiseless training signs,
\begin{equation}\label{eq:draw}
 \E[\|G\|_K^2\mid\text{global margins and training signs}]\ge M.
\end{equation}
Indeed write the whitened Gaussian as $K^{-1/2}G=RU$, where the radius $R$ and direction $U$ are independent and $\E R^2=M$. For each direction satisfying the signs, the global margins impose only a lower cutoff on $R$. Truncating a nonnegative random variable above a lower cutoff increases its mean square. Averaging over the conditional directions proves~\eqref{eq:draw}. Thus draw-energy estimates must specify a dimension baseline or a centered or projected quantity.

The obstruction proved here concerns amplitude and energy. Predictive correctness with probability strictly greater than $2/3$ is a separate quantity: a large posterior mean at the observed input does not imply confident correct prediction on unobserved inputs. No sample-complexity conclusion is inferred from this note.
\end{document}
